% GENERATED FILE. Edit canonical lessons, then run npm run generate:books.
% canonical-source-hash: fnv1a64:388a6950ab136cc1
% canonical-lessons: GU-C06-numbers-1-5, GU-C06-number-histories

\chapter{Numbers One to Five}
\label{ch:numbers-one-to-five}

This chapter is generated from the canonical micro-lessons used by Language Ladder.
Each section stays independently resumable and preserves the authored prerequisite order.

\section[ek be traṇ chār pā̃ch]{\gu{એક}, \gu{બે}, \gu{ત્રણ}, \gu{ચાર}, \gu{પાંચ} — one to five}

\label{lesson:GU-C06-numbers-1-5}



\begin{quote}
\textbf{Warm-up.} {[}PAUSE 2s{]} Three of these will look familiar from any other Indian language. Two won't — and both have a reason.
\end{quote}

\subsection*{You'll want to know: the five}
\noindent
\begin{tabularx}{\linewidth}{@{}>{\raggedright\arraybackslash}X>{\raggedright\arraybackslash}X>{\raggedright\arraybackslash}X@{}}
\toprule
\textbf{} & \textbf{Gujarati} & \textbf{said} \\
\midrule
1 & \textbf{\gu{એક}} & \emph{ek} \\
2 & \textbf{\gu{બે}} & \emph{be} \\
3 & \textbf{\gu{ત્રણ}} & \emph{traṇ} \\
4 & \textbf{\gu{ચાર}} & \emph{chār} \\
5 & \textbf{\gu{પાંચ}} & \emph{pā̃ch} \\
\bottomrule
\end{tabularx}

Notice the script: Gujarati is Devanagari \textbf{without the top line}. Same letters, same system — the shirorekhā simply isn't drawn. If you've done the Hindi writing track you can feel the relationship immediately.

\begin{grammarlens}[title={What you've built: two surprises and a headless script}]
Most neighbouring languages have a \emph{d} in “two,” but Gujarati says \textbf{\gu{બે}} \emph{be}. Hindi says \emph{tīn} for “three,” but Gujarati says \textbf{\gu{ત્રણ}} \emph{traṇ}, with an \emph{r}.

Do not carry both histories while the count is still new. Make the five forms automatic first; the next prerequisite-ordered lesson explains why \emph{be} and \emph{traṇ} took those shapes.
\end{grammarlens}

\subsection*{Guided Practice}
{[}PAUSE 1s{]}

\begin{itemize}
\raggedright
  \item {[}YOU SAY: "ek, be, traṇ, chār, pā̃ch"{]}
  \item {[}YOU SAY: the odd two — Gujarati \textbf{be} and \textbf{traṇ}{]}
  \item {[}YOU SAY: the script clue — Devanagari shapes without a top line{]}
\end{itemize}

\begin{tcolorbox}[breakable,colback=teal!4,colframe=teal!35!black,title={Wrap-up recall}]
{[}PAUSE 3s{]} Count to five in Gujarati. (\emph{Ek, be, traṇ, chār, pā̃ch}.) Why does \textbf{\gu{બે}} look unlike its neighbours? (It begins with \emph{b}, not \emph{d}.) What does \textbf{\gu{ત્રણ}} contain that Hindi \emph{tīn} does not? (\emph{r}.) What is visibly missing from the Gujarati script? (The top line: the \emph{shirorekhā} is not drawn.) Next: trace the two surprising number histories.
\end{tcolorbox}
\section[be · traṇ]{Two number histories: \emph{be} and \emph{traṇ}}

\label{lesson:GU-C06-number-histories}



\begin{quote}
\textbf{Warm-up.} {[}PAUSE 2s{]} Count once: \emph{ek, be, traṇ, chār, pā̃ch}. Which two forms look least like Hindi or Marathi? (\emph{Be} and \emph{traṇ}.)
\end{quote}

\begin{cousinweb}
Sanskrit had gendered forms of “two,” and its daughters selected different parts of that paradigm:

\noindent
\begin{tabularx}{\linewidth}{@{}>{\raggedright\arraybackslash}X>{\raggedright\arraybackslash}X>{\raggedright\arraybackslash}X@{}}
\toprule
\textbf{language} & \textbf{“two”} & \textbf{route} \\
\midrule
Sanskrit & \emph{dváu} (masc.) · \emph{dvé} (fem. and neut.) & — \\
Hindi & \emph{do} & masculine side \\
Marathi & \emph{don} & masculine side, plus an analogical \emph{-n} \\
Bengali & \emph{dui} & disyllabic Prakrit \emph{duve} \\
\textbf{Gujarati} & \emph{\textbf{be}} & feminine/neuter \emph{dvé} side \\
\bottomrule
\end{tabularx}

Gujarati \textbf{\gu{બે}} is not a corruption of \emph{do}. It is a \textbf{different inheritance} from the same Sanskrit paradigm.

The initial cluster changed too. In \emph{dv-}, the dental \emph{d} adopted the \textbf{labial place} of the following \emph{v}: \emph{dv} $\to$ \emph{bb} $\to$ \emph{b}. The first consonant was remade in the shape of its neighbour, then the double simplified.
\end{cousinweb}

\begin{cousinweb}
\noindent
\begin{tabularx}{\linewidth}{@{}>{\raggedright\arraybackslash}X>{\raggedright\arraybackslash}X@{}}
\toprule
\textbf{stage} & \textbf{“three”} \\
\midrule
Sanskrit neuter & \emph{trī́ṇi} \\
Prakrit & \emph{tiṇṇi}: \emph{r} already gone \\
Hindi & \emph{tīn} \\
\textbf{Gujarati} & \emph{\textbf{traṇ}}: \emph{r} present again \\
\bottomrule
\end{tabularx}

The shared \textbf{ṇ} shows that Hindi and Gujarati both descend from neuter \emph{trī́ṇi} through Prakrit \emph{tiṇṇi}. By that stage, the \emph{r} was gone and its weight had moved into doubled \emph{ṇṇ}. Hindi later simplified the double and lengthened the vowel.

Gujarati's \emph{tr-} is generally treated as \textbf{restored from Sanskrit}, not carried through unbroken. Sanskrit remained a literary language and continued to reshape its descendants. \emph{Traṇ} looks older than \emph{tīn}, but it is closer because speakers put the \emph{r} back.
\end{cousinweb}

\subsection*{Guided Practice}
{[}PAUSE 1s{]}

\begin{itemize}
\raggedright
  \item {[}YOU SAY: Gujarati's selected form — feminine/neuter \emph{dvé}{]}
  \item {[}YOU SAY: the sound path — \emph{dv $\to$ bb $\to$ b}{]}
  \item {[}YOU SAY: the \emph{r} path — \emph{trī́ṇi $\to$ tiṇṇi $\to$ traṇ}{]}
  \item {[}YOU SAY: inherited continuously or restored — restored{]}
\end{itemize}

\begin{tcolorbox}[breakable,colback=teal!4,colframe=teal!35!black,title={Wrap-up recall}]
{[}PAUSE 3s{]} Why \emph{be}, not Hindi \emph{do}? (Different Sanskrit forms.) What changed in \emph{dv $\to$ b}? (\emph{D} became labial beside \emph{v}; the double simplified.) Did \emph{traṇ} keep \emph{r}? (No: Sanskrit restored Prakrit's lost \emph{r}.) How can a newer form look older? (Learned restoration.)
\end{tcolorbox}
